% Revised working copy of F024. The original ZIP remains unchanged.
\documentclass[12pt,a4paper]{article}
% Shared print layout. The web's provisional design and MathJax are independent.
\usepackage[utf8]{inputenc}
\usepackage[T1]{fontenc}
\usepackage[spanish,es-nodecimaldot]{babel}
\usepackage{lmodern}
\usepackage{amsmath,amssymb}
\usepackage{graphicx}
\usepackage{geometry}
\geometry{a4paper,left=20mm,right=20mm,top=25mm,bottom=22mm,headheight=16pt,headsep=8mm}
\usepackage{fancyhdr,lastpage}
\usepackage{enumitem}
\usepackage{xcolor}
\usepackage{pgfplots}
\pgfplotsset{compat=1.18}
\usepgfplotslibrary{fillbetween}
\usepackage[hidelinks,pdfencoding=auto]{hyperref}
\hypersetup{pdflang=es-ES,pdfauthor={},pdfsubject={Matemáticas Aplicadas a las Ciencias Sociales II}}
\pagestyle{fancy}
\fancyhf{}
\fancyhead[L]{\small Matemáticas Aplicadas a las Ciencias Sociales II}
\fancyhead[R]{\small 2.\textsuperscript{o} Bachillerato}
\fancyfoot[L]{\footnotesize Recuperación T1 | Edición revisada 1.0}
\fancyfoot[R]{\footnotesize Página \thepage\ de \pageref{LastPage}}
\renewcommand{\headrulewidth}{0.4pt}
\setlength{\parindent}{0pt}
\setlength{\parskip}{0.5em}
\setlist{itemsep=0.3em,topsep=0.3em}
\newcommand{\apartado}[1]{\subsection*{#1}}
\pgfplotsset{
  revision/.style={
    width=0.94\linewidth,height=7.5cm,
    axis lines=middle,grid=major,grid style={gray!18},
    tick label style={font=\small},label style={font=\small},
    xlabel={$x$},ylabel={$y$},samples=160,
    legend style={at={(0.5,-0.16)},anchor=north,draw=none,font=\small},
    clip=true
  }
}

\hypersetup{pdftitle={Recuperación T1 - Examen revisado}}
\begin{document}
\begin{minipage}[c]{0.17\linewidth}
  \includegraphics[width=\linewidth]{logo-al-andalus.png}
\end{minipage}\hfill
\begin{minipage}[c]{0.78\linewidth}
  {\Large\bfseries Examen de recuperación T1}\\[0.5em]
  Matemáticas Aplicadas a las Ciencias Sociales II\\
  \textbf{Duración: 55 minutos \quad Puntuación total: 10 puntos}
\end{minipage}

\vspace{1em}
\textbf{Apellidos y nombre:} \hrulefill

\textbf{Fecha:} \rule{2cm}{0.4pt} de \rule{4cm}{0.4pt} de \rule{1.7cm}{0.4pt}

\vspace{0.5em}
\fbox{\parbox{\dimexpr\linewidth-2\fboxsep-2\fboxrule\relax}{
Lee atentamente todas las preguntas y respóndelas de forma razonada.
Redacta lo necesario para que tus razonamientos sean claros.
Numera las páginas de tu respuesta. Se permite el uso de calculadora.
Tienes un máximo de \textbf{55 minutos} para terminar.
}}

\section*{Ejercicio 1 \hfill 4 puntos}
\textit{Tiempo orientativo: 15 minutos.}

Se considera la función
\[
f(x)=
\begin{cases}
ax-3x^2,& x\leq 1,\\
2x^2+b,& x>1.
\end{cases}
\]
\begin{enumerate}[label=\textbf{(\alph*)},leftmargin=*,itemsep=1em]
  \item \textbf{(2 puntos)} Calcula los valores de \(a\) y \(b\) para que
  \(f\) sea derivable en \(x=1\).
  \item \textbf{(2 puntos)} Para \(a=3\) y \(b=-2\), halla la ecuación de la
  \textbf{recta tangente} a la gráfica de \(f\) en el punto de abscisa \(x=2\).
\end{enumerate}

\section*{Ejercicio 2 \hfill 6 puntos}
\textit{Tiempo orientativo: 40 minutos.}

Dada la función
\[
g(x)=\frac{1}{x+3},
\]
\begin{enumerate}[label=\textbf{(\alph*)},leftmargin=*,itemsep=1em]
  \item \textbf{(4 puntos)} Estudia la continuidad, las asíntotas,
  la monotonía y la curvatura de \(g\), y represéntala gráficamente.
  \item \textbf{(2 puntos)} Calcula el área comprendida entre la gráfica
  de \(g\), el eje de abscisas y las rectas \(x=2\) y \(x=4\).
\end{enumerate}
\end{document}
